% Klumpen, ei lignelse om form
% Djupt prosaisk parabel som fangar heile den akademiske historia til designfaget,
% berre om form (klumpen/gruten). Utan namn, signatur eller ornament.
% Bygg: xelatex klumpen.tex
\documentclass[11pt]{article}

\usepackage[paperwidth=148mm,paperheight=210mm,top=20mm,bottom=20mm,left=22mm,right=22mm]{geometry}
\usepackage{fontspec}
\usepackage{polyglossia}\setdefaultlanguage[variant=nynorsk]{norwegian}
\IfFontExistsTF{Source Serif 4}{\setmainfont{Source Serif 4}}{%
  \IfFontExistsTF{EB Garamond}{\setmainfont{EB Garamond}[AutoFakeBold=2.5]}{%
  \setmainfont{TeX Gyre Pagella}}}
\usepackage{microtype}\usepackage{xcolor}

\setlength{\parindent}{0pt}\setlength{\parskip}{8pt}\linespread{1.12}
\pagestyle{empty}\frenchspacing

\begin{document}

{\footnotesize\addfontfeature{LetterSpace=22}\color{gray}ei lignelse om form}

\vspace{3mm}
{\fontsize{22}{25}\selectfont\bfseries Klumpen}

\vspace{2mm}\noindent{\color{gray}\rule{14mm}{0.4pt}}

\vspace{5mm}

{\large Det var ein klump.}

Før kvar form, før kvart namn, låg han der, og han var ikkje ingenting. Han var alt som kunne bli og enno ikkje var. Held du han i handa, heldt du eit heilt rom, alle formene klumpen kunne gå inn i, og dei var mest alle saman tomme, for det meste som kan bli, blir aldri. Slik byrja faget, med ein klump og eit val ingen heilt forstod at dei tok.

Den fyrste handa kjende klumpen utan ord. Ho tenkte i leira, ho visste og laga i same rørsla, og ho spurde aldri kvifor nett denne forma og ikkje ei anna, for svaret sat i handa og ikkje i munnen. Det var ein god kunnskap. Men det var ein kunnskap som døydde med den som bar han, for handa kunne vise, og ho kunne ikkje seie.

Så vart handa delt. Ein teikna, ein annan pressa, ein tredje selde, og klumpen gjekk frå hand til hand til ingen av dei lenger tok på det forma gjorde når ho forlét dei. Skaden klumpen bar, bar han til ein annan stad, ein annan kropp, eit anna år, og i bøkene stod berre det som vart selt. Det som vart øydelagt, stod ingen stad.

Så kom skulane, ein etter ein, kvar med si fane, og kvar stilte seg framfor klumpen og lova det same, at no skulle dei finne grunnen, no skulle dei seie kvifor denne forma og ikkje ei anna. Den eine kalla forma ærleg. Den andre kalla henne funksjonell. Den tredje kalla henne god. Men eit namn er ingen grunn, og den som sa at forma følgjer det ho gjer, sa ei setning som høver på kvart utfall, og ei setning som høver på alt, seier ingenting.

Ein gong, berre ein gong, kom det nokon og la eit mål ned ved sida av klumpen, ein måte å seie kvar i rommet forma låg, kva trykk som drog henne dit, kva som ville hende om eitt trykk gav etter. Faget tok prestisjen i ordet og lét målet liggje på bordet. Det var lettare å halde fram med å døype.

Og då nokon spurde kvifor, svara dei til slutt det vakraste eit fag har funne for å sleppe unna, at klumpen er for vrang til å målast, at kunnskapen er for taus til å seiast, at dette lærer du ved benken, med tida. Dei gjorde det dei ikkje kunne grunngje om til ein dygd. Og det einaste dei greidde å sende ut or huset, var ein tom framgangsmåte ingen trong kunnskap for å bruke, og han slo rot der det fanst ei ramme å feste han i, langt frå klumpen.

Men heile tida heldt klumpen fram med å gå ut i verda, og han verka. Han verka langs den eine aksen dei bad om, og han svikta langs alle dei andre. Koppen som tek tolv sekund å laga og fire hundre år å brytast ned. Fjella av ting vi eig eit år og så er framande for. Maskina som et merksemda til borna. Forma sviktar der ingen såg, og svikten er leseleg, slik eit avtrykk i stein ber dyret som låg der. Du kan lese kor langt omhugen rakk, rett av forma.

Hundre år gjekk, og enno fanst det ikkje eit felles ord for klumpen, ikkje eit mål, ikkje eit råd som kunne seie nei. Dei bar fråvêret som ein lagnad og kalla det faget si eiga natur.

Men klumpen var aldri grunnlaus. Han var alltid ein posisjon, eitt punkt i eit rom av alt han kunne vore, og det tomme i rommet var ikkje det umoglege, det var det moglege ingen nådde. Det som flytta han, var aldri funksjonen og aldri hundreåret, det var alt som drog i han på ein gong. Og det eine draget dei aldri gav namn og aldri målte, var omhugen, gapet mellom forma dei laga og forma verda trong.

Å lage godt er å vite kva veg ein skal dra. Å døme godt er å seie kva akse, og kva det kosta. Når du ikkje kan seie kvifor, har du ikkje dømt. Du har berre gjeve klumpen eit namn.

\medskip

Så mål han. Sjå til botnen, til gruten, til grunnen som låg der heile tida.

{\itshape Forma lyg ikkje, og fråvêret av henne lyg ikkje det heller.}

\end{document}
